% Formal model of evidence-gated completion (Sec. IV-C). Code mapping: artifacts/07_claims_ledger.md.
\subsection{Evidence-Gated Verification Model}
\label{sec:formal}

We model an episode as a sequence of browser states $s_t$ with observations $o_t=\Omega(s_t)$ (URL, title, element manifest, screenshot). The planner maps the goal $g$, $o_t$ and retrieved context to an action $a_t$; the executor applies it and returns a dispatch flag $\alpha_t\in\{0,1\}$, which is $1$ when the browser call raised no exception and, for text entry, the value read back matched. Some fallback paths (a mouse-and-keyboard retry of text entry, vision-coordinate clicks) return $\alpha_t=1$ without such a check.

\paragraph{Blocking predicate} $\beta(s)=1$ when the page is a bot check: the URL contains a challenge path, the title contains \emph{captcha}, \emph{robot} or \emph{sorry\ldots}, the body contains one of 12 challenge phrases, or a challenge frame is present. This check runs before $\Phi$ and ends the episode as \textsc{Blocked}.

\paragraph{Step and task predicates} The decomposer produces a plan $P=(g_1,\dots,g_n)$. The per-step predicate is currently weak, $\phi_i := \alpha_t$: a step is marked complete when its action dispatches. Completion is decided by $\Phi(g,s_{t+1},\hat{y},P)$, evaluated on the live page after the action. With $\hat{y}$ the answer,
\begin{equation}
\Phi = \neg E(s)\;\wedge\;\neg H(g,s,P)\;\wedge\;T(g,s)\;\wedge\;X(g,\hat{y})\;\wedge\;R(P,\hat{y}),
\label{eq:phi}
\end{equation}
where $E$ holds for a browser-error URL or a title containing one of four HTTP error phrases (404, 500, 502, 503); $H$ holds when the request contains a search or extraction keyword, or the plan has more than one step, and the URL is the bare home page of one of four search engines; $T$ requires that the literal text extracted from $g$ by a pattern (``enter exactly \ldots'', ``type \ldots'') appears, as a substring, in an input, text area or editable element; $X$, instantiated by extraction keywords, requires a non-empty structured extraction or an answer longer than 40 characters that does not begin with a navigation verb; and $R$ requires that no plan step is pending unless an answer exists. $H$, $T$, $X$ and $R$ are instantiated only when their triggers are present; for a request that triggers none of them, $\Phi$ reduces to $\neg E$.

\paragraph{When $\Phi$ is evaluated} $\Phi$ is evaluated when the planner emits \texttt{complete}, when the plan index has reached $n$, or when a requested text entry has dispatched. For a one-step plan, the default for short requests, this is after every successful action, so $\Phi$ also acts as the termination rule. The answer is $\hat{y} = y_{\mathrm{ext}}$ if an extraction has produced one, otherwise the planner's \texttt{reasoning} string for the last action, which makes $X$ and $R$ easy to satisfy (Section~\ref{sec:results}).

\paragraph{Completion rule} The agent reports success only if $\Phi=1$; each evaluation is written to the evidence store with its expected and observed values. A passing evaluation can still yield \textsc{Waiting approval}, and the terminal node downgrades a success to \textsc{Failed} if the intent named a site but no navigation was verified. A rejected evaluation does not end the episode: control returns to observation and planning until a counter of intent, decomposition and planning calls reaches $K=15$.

\paragraph{Consequence} Dispatch success does not imply reported success, $\alpha_t = 1 \not\Rightarrow \Phi=1$; a successful \texttt{navigate} to a search engine's home page for a search request is an example ($H=1$). The converse guarantee, that $\Phi=1$ implies the goal is met, does not hold, and Section~\ref{sec:results} measures how often it fails.

\paragraph{False completion rate} Let $\mathcal{T}_c$ be the episodes reported complete and $\Phi^{\star}$ an independent, task-specific oracle written by the experimenter. Then
\begin{equation}
\mathrm{FCR} = \frac{\lvert\{\tau\in\mathcal{T}_c : \Phi^{\star}(\tau)=0\}\rvert}{\lvert\mathcal{T}_c\rvert}.
\label{eq:fcr}
\end{equation}
FCR is a precision-type measure; an agent that never completes has no false completions. It must therefore be reported together with task success and the false-rejection rate, the fraction of goal-satisfying episodes not reported complete.
